\scanpage{136}
\mspage{85}

\deleted{de $X$ dont l'ens. d'objets est $\mathrm{Ob}\,U\cup\{t\}$, où
$t$ \deleted{\unsure{$\in\dot U$|$\in\overline U$} \unsure{th.}} est un élément donné de $\dot U$. \deleted{Alors}
\unsure{(}Donc $U$ est un ouvert dans $X_1$, \deleted{\unsure{dense dans $X_1$},}
dont l'adh. dans $X_1$ est égale à $X_1$, et le
complémentaire est \deleted{la catégorie \unsure{ponc}} réduit à $t$.
Alors $U\xrightarrow{\unsure{i_1|i}}X_1$ est dans $\mathcal W$ sss $U_{/t}$ est
$\mathcal W$-asphérique. Même résultat pour la
catégorie $X'_1(t)$ définie par le foncteur
$e\to T$ d'image $t$ (comme dans b'').}

\deleted{Dém. Par construction, la catégorie $U_{/t}$
au sens du couple $(X,U)$ est isomorphe
à celle au sens du couple $(X_1(t),U)$, ou
du couple $(X'_1(t),U)$. Donc on est ramené
à ceci :}

\uline{Corollaire \deleted{\unsure{2}}} \inserted{(Moyennant $\mathrm W(7\ \mathrm{bis})$)}
Soit $X$ dans $\mathrm{Cat}$, $U$ un ouvert de $X$
tel que $X\setminus U=$\deleted{$\unsure{\overline{\{t\}}|\{t\}}$, et $X=\overline U$, \unsure{i.e.} $t\in\overline U$.}%
\inserted{soit la catégorie finale.}
Alors l'inclusion $U\overset{i}{\unsure{-|\to}}X$ est dans $\mathcal W$ sss
elle est $\mathcal W$-asphérique, i.e. sss $U_{/t}$ est $\mathcal W$-asphérique. (où $t$ est l'unique élément de $T$).

Nous démontrerons ultérieurement ce
corollaire\textsuperscript{*}.
% slanted
\margin{cf cor 4 de la prop.~15, page 149}
La prop.~7 ne nous servira pas
d'ici là -- elle n'aura qu'un rôle
purement \unsure{esthétique} à jouer, pour justifier
notre propos délibéré de chercher des
conditions sur $i:U\to X$ \deleted{\unsure{qui soient}}\inserted{\unsure{qui assurent|exprimant}}
que $U\to\overline U$ est dans $\mathcal W$ et le reste, \deleted{\unsure{par}}
quand on \deleted{fait bouger}\inserted{remplace} $(X,U)$ par d'autres
couples $(X_1,U_1\unsure{.})$ qui s'en déduisent par
certaines opérations \ldots) qui

% archive sheet number: 136
